\begin{minipage}{0.60\linewidth}
    Un pendule est constitué d'une petite boule métallique, de masse
    $m=\qty{80}{\g}$, suspendue à un fil inextensible, de masse considérée comme
    négligeable et de longueur $\ell=\qty{1.00}{\m}$.
    
    Le fil est accroché en un point fixe $O$ et les mouvements du pendule
    s'effectuent dans un plan vertical. Le fil du pendule étant initialement
    vertical, on l'écarte de cette position d'un angle $\theta_{m}=\ang{45}$.
    
    Puis, fil tendu, on le lâche sans vitesse (position~$A$).
\end{minipage}
\hfill
\begin{minipage}{0.38\linewidth}
\begin{figure}[H]
    \centering
    \begin{tikzpicture}
        \def\tho{45}
        \newdimen\el \el=4cm
        \node[draw,thick,pattern=north east lines,
              minimum width=1cm,minimum height=5mm] (S) {};
        \coordinate (O) at (S.south);
        \path[draw,fill=DarkTurquoise]
            (O) -- ([yshift=-1cm]O) arc (270:270+\tho:1cm) --
            cycle;
        \node at ($(O)+(270+\tho/2:1.4)$) {$\theta_{m}$};
        \node[circle,fill,inner sep=1.1pt,
              label={[anchor=north east,inner sep=1.8pt]:$O$}]
              at (O) {};
        \draw[densely dashed]
            (O) +(270:\el) arc (270:270+\tho:\el);
        \foreach \p/\a in {A/\tho,B/0}{%
            \node[circle,ball color=black!60,inner sep=0pt,
                  minimum size=4mm,
                  label={[inner sep=1.6pt]-90+\a:$\p$}]
                  (\p) at ($(O)+(270+\a:\el)$) {};
            \fill (\p.center) circle (1pt);
            \draw[thick] (O) -- (\p);
        }
        \node[anchor=180+\tho,inner sep=1.4pt]
            at ($(O)!.5!(A.center)$) {$\ell$};
        \draw[densely dashed] (A) -- (A -| O) coordinate (H)
            node[anchor=east] {$H$};
        \draw[thick] (H) +(0,-6pt) -| +(6pt,0);
        \node[right=2mm] at (A) {$(v=0)$};
        \draw[densely dashed] ([xshift=-1.2cm]B.center) -- ++(4,0)
            node[signal,draw,fill=lightgray,anchor=west,solid,
            signal to=west] {$E_{pp}=0$};
        \draw[->,thick,blue] (B.center) -- ++(-0.8,0)
            node[pos=0.8,below] {$\vec{v}^{\prime}$};
    \end{tikzpicture}
\end{figure}
\end{minipage}


\begin{enumerate}
    \item Décrire le mouvement de la boule.

          Justifier la conservation de la somme $E_{c}+E_{pp}$ pour la boule du
          pendule.
    
          Quelle est la transformation d'énergie qui s'effectue au cours du
          mouvement~?
    \item Déterminer la valeur $v^{\prime}$ de la vitesse de la boule
          lorsqu'elle passe par la position verticale (position~$B$).
    
          Justifier le calcul.
    
          Faire le calcul littéral et procéder à l'application numérique.
\end{enumerate}
